% Generated draft: requires mathematical and source review.
\documentclass[11pt]{article}
\usepackage[margin=25mm]{geometry}
\usepackage{fontspec}
\setmainfont{Latin Modern Roman}
\usepackage{amsmath,amssymb,mathtools,unicode-math}
\setmathfont{Latin Modern Math}
\usepackage{enumitem,needspace,xurl,hyperref}
\setlength{\parindent}{0pt}
\setlength{\parskip}{5pt}
\setlength{\emergencystretch}{3em}
\newcommand{\N}{\mathbb{N}}
\newcommand{\Z}{\mathbb{Z}}
\newcommand{\Q}{\mathbb{Q}}
\newcommand{\R}{\mathbb{R}}
\newcommand{\C}{\mathbb{C}}
\newcommand{\F}{\mathbb{F}}
\newcommand{\A}{\mathbb{A}}
\newcommand{\PP}{\mathbb{P}}
\newcommand{\E}{\mathbb{E}}
\newcommand{\eps}{\varepsilon}
\newcommand{\dd}{\,\mathrm{d}}
\begin{document}
\section*{UnsolvedMath ID: 30005538}
\textbf{Source problem code:} catalogue entry
\subsection{Definitions and mathematical statement}
Let \(L\to X\) be a holomorphic line bundle with a smooth Hermitian metric \(h^L\), and let \(R^L\) denote its Chern curvature. Semipositivity means that the real \((1,1)\)-form \(\sqrt{-1}R^L\) is nonnegative. Let \((F,h^F)\to X\) be a Hermitian holomorphic vector bundle. Once a volume form on \(X\) has been specified, the metrics define an \(L^2\) inner product on sections of \(L^{\otimes p}\otimes F\), for positive integers \(p\). In the compact setting, let \(P_p\) be the orthogonal projection onto \(H^0(X,L^{\otimes p}\otimes F)\), and write \(P_p(x,x)\in\operatorname{End}(F_x)\) for its diagonal Bergman kernel, using the natural identification of endomorphisms of the one-dimensional factor \(L_x^{\otimes p}\).

\textbf{Requested higher-dimensional extension.} Establish a full pointwise asymptotic expansion of \(P_p(x,x)\) as \(p\to\infty\) at curvature-vanishing points for semipositive line bundles when \(\dim_{\C}X>1\). Here a full pointwise expansion means an asymptotic description to arbitrarily high order at each fixed point under consideration. The supplied record does not specify the higher-dimensional hypotheses, the permitted asymptotic powers or other terms, or the remainder estimates. Thus this is a research target requiring further formulation, rather than a completely specified universally quantified conjecture. Compactness is inherited from the motivating surface theorem below; it is not explicitly imposed by the supplied higher-dimensional statement. The volume form and any additional geometric assumptions must also be specified before fixing the kernel normalization.

\textbf{Supplied motivating surface result.} The excerpt states that if \(X\) is a compact Riemann surface, \((L,h^L)\to X\) is semipositive with curvature vanishing to finite order at every point, and \((F,h^F)\to X\) is a Hermitian holomorphic vector bundle, then, at each fixed \(x\in X\),
\[
P_p(x,x)=\sum_{j=0}^{N}c_j(x)p^{(2-2j)/r_x}
 +O\!\left(p^{-2N/r_x}\right)
\qquad(p\to\infty)
\]
for every truncation index \(N\). The excerpt describes the \(c_j\) as sections of \(\operatorname{End}(F)\), states that \(c_0(x)>0\), and says that this leading coefficient is determined by the Bergman kernel of a model Kodaira Laplacian on \(T_xX\). Positivity here is positivity of the Hermitian endomorphism. The parameter \(r_x\), including its precise relation to curvature vanishing order, is not defined in the supplied excerpt. This formula is recorded as the supplied surface result, not asserted as the desired higher-dimensional formula.

The excerpt also supplies a surface spectral-gap assertion under the same hypotheses, with positive constants \(c_1,c_2\) and lower threshold \(c_1p^{2/r}-c_2\), but does not define its global parameter \(r\) or clearly identify the operator in the extracted text. It reports that a corresponding spectral gap fails in higher dimensions, citing Donnelly's counterexample to Siu's eigenvalue conjecture. Consequently, the surface spectral-gap assertion cannot be silently adopted as a hypothesis-free higher-dimensional tool.

The higher-dimensional target leaves unresolved whether curvature vanishing means \(R^L(x)=0\) as a tensor or includes points where a nonzero semipositive curvature form has deficient rank. It also leaves unresolved whether finite-order vanishing is required at every point, only at the point of expansion, or in specified directions, and whether twisting by arbitrary \(F\) is intended.
\subsection{Short English statement}
The problem asks for arbitrarily detailed pointwise Bergman-kernel asymptotics at curvature-vanishing points in complex dimension greater than one. A supplied theorem provides the motivating expansion on compact Riemann surfaces with finite-order curvature vanishing. The higher-dimensional geometric assumptions and the form of the expansion remain unspecified, so this draft records the target and its formulation gaps without asserting a new theorem.
\subsection{Sources}
\begin{itemize}
\item Local registry: UnsolvedMath permanent ID \(30005538\), source code \texttt{catalogue entry}, title \emph{Bergman-Kernel Asymptotics at Curvature-Vanishing Points}; supplied complete source record. No individual proposer is supplied; the proposed year is 2023.
\item Supplied report citation: \emph{Hypoelliptic Operators in Geometry} (2023), Oberwolfach Reports, pp.\ 1323--1396, DOI \texttt{10.4171/owr/2023/24}, \url{https://doi.org/10.4171/owr/2023/24}. The record attributes the problem to workshop-report material.
\item Supplied literature source, labeled ``OWR 2023 24'': \url{https://ems.press/content/serial-article-files/47017}.
\item Supplied literature source, labeled ``Marinescu Savale'': \url{https://doi.org/10.1007/s00208-023-02750-3}. No article title or fuller bibliographic record is supplied.
\item The supplied excerpt labels its surface results Corollary 3 and Theorem 4, both attributed there to reference [13], and points to [9, Problem 4.8] for the broader expansion problem. It also mentions [9, Theorem 1.5.5] and methods from [5, 9, 10]. The numbered bibliography is not included, so these labels cannot be fully resolved from the record.
\item The excerpt attributes an earlier asymptotic estimate for semipositive line bundles to Berndtsson; an expansion on the positive part and away from the augmented base locus, assuming ampleness, to Berman; and an expansion on the positive part with canonical-bundle twisting \(F=K_X\) to Hsiao--Marinescu. It attributes the higher-dimensional spectral-gap counterexample to Donnelly, reference [6]. These are supplied contextual claims without complete individual citations.
\item Literature assessment dated \textbf{2026-08-22}: the record labels the problem open and says that the Riemann surface case is proved while the corresponding higher-dimensional semipositive expansion remains unresolved. This is supplied dated curation, not an independently verified literature assessment. The record expressly requests independent verification of time-sensitive status and cited claims.
\item The registry labels the statement \texttt{corrected\_verified} and says that it was reconstructed and checked against the official report. These are supplied provenance claims; no independent source inspection or mathematical verification is asserted here.
\end{itemize}
\end{document}
